\documentclass[10pt,a4paper]{amsart}
\usepackage[margin=19mm]{geometry}
\usepackage{amsmath,amssymb,mathtools}
\usepackage[scr=rsfs,cal=euler]{mathalfa}
\usepackage{xfrac}
\usepackage[unicode,hidelinks,bookmarks=false]{hyperref}
\newcommand{\ie}{i.e.\ }
\newcommand{\cf}{cf.\ }
\newcommand{\resp}{respectively\ }
\newcommand{\Subsection}{\subsection}
\providecommand{\nobreakdash}{\nobreak\hbox{-}}
\setlength{\emergencystretch}{2em}
\multlinegap0pt
\usepackage{bm}
\usepackage{bbm}
\usepackage{dsfont}
\usepackage{xspace}
\usepackage{cancel}
\usepackage{mathtools}
\usepackage[shortlabels]{enumitem}
\usepackage{amsthm}
\makeatletter
\@ifundefined{theorem}{\newtheorem{theorem}{Theorem}}{}
\@ifundefined{proposition}{\newtheorem{proposition}[theorem]{Proposition}}{}
\makeatother
\renewcommand{\Re}{\mathbb{R}}
\title[Submodular and strongly submodular functions and diversities]{Submodular and strongly submodular functions and diversities}
\author{David Bryant, Paul F. Tupper}
\date{Source: arXiv:2608.11468v1 (CC BY 4.0)}
\begin{document}
\maketitle

\noindent\emph{Source and licence.} Submodular and strongly submodular functions and diversities. Version arXiv:2608.11468v1 (CC BY 4.0), distributed under the Creative Commons Attribution 4.0 International licence, as recorded in the arXiv licence metadata for this exact version and shown on the abstract page https://arxiv.org/abs/2608.11468v1. Full rights notice: licenses/arxiv-2608.11468-CC-BY-4.0.txt.\par\smallskip

\noindent\textit{Editorial scope.} Counted unit: the Theorem whose source label is \texttt{thm:submodular\_match} in the article (source0.tex lines 668--670, source label \texttt{thm:submodular\_match}), delivered together with the article's own complete original proof of it (source0.tex lines 671--712, one \texttt{proof} environment of 42 lines). No mathematics is written, repaired or re-derived: statement and proof are the article's own text line for line. Required context retained uncounted: prop:basic\_strong (lines 330--349). Every definition, display and label of those lines is retained unchanged. Omitted from this package and not redistributed: 1--329, 350--667, 713--942. Nothing inside a retained excerpt is omitted or altered: every retained body is delivered exactly as the article's lines are written, so no omission receipt of any kind accompanies this package. Citations: the retained excerpts carry no citation command at all, so no bibliography entry of the article is transcribed and no reference is redistributed. Reference adaptation: none was needed and none was made; every reference inside the delivered unit resolves inside it, so no unresolved reference or citation is produced. Printed slips kept literal: no printed word, symbol, hypothesis or number of the counted statement or of its proof was altered. Layout and class: the article's own class is \texttt{article} with packages including graphicx, enumitem, amsthm, color, times, authblk, hyperref, url, setspace, amsmath, amssymb, relsize, mathtools; the wrapper is the lane's pinned \texttt{amsart} editorial wrapper at 10pt on A4, which restates the theorem-like environments the retained text needs and adds the article's own macro definitions that the retained text needs (\texttt{\textbackslash Re}), with extra wrapper packages (bm, bbm, dsfont, xspace, cancel, mathtools, enumitem). The delivered numbering is the unit's own. Assets: no retained span contains a figure environment, a graphics-inclusion command, a PDF-page inclusion, a table or a TikZ picture; no image file is copied into this package, the package declares no asset dependency and no image file is redistributed with it. Licence: version arXiv:2608.11468v1 of this article is distributed under the Creative Commons Attribution 4.0 International licence, as recorded in the arXiv licence metadata for this exact version and shown on the abstract page https://arxiv.org/abs/2608.11468v1. A full rights notice is delivered at licenses/arxiv-2608.11468-CC-BY-4.0.txt. Characters: every retained excerpt is pure ASCII, so the delivered engine, XeLaTeX with its default Latin Modern text font, reports no missing character.
\begin{proposition} \label{prop:basic_strong}
Let $X$ be a finite set.
\begin{enumerate}[label=(\roman*)]
\item For all $\alpha \in \Re$ the constant function $f(A) = \alpha$ is strongly submodular.
\item If $f_1,f_2$ are strongly submodular and $\alpha \geq 0$ then $\alpha f_1$ and $f_1 + f_2$ are strongly submodular.
\item Every strongly submodular function $f:2^X \rightarrow \Re$  is submodular.
\item Every strongly submodular function $f:2^X \rightarrow \Re$ is non-decreasing. %pg 73
\item If $\Gamma:2^X \rightarrow 2^Y$ satisfies $\Gamma(A \cup B) = \Gamma(A) \cup \Gamma(B)$ for all $A,B \in 2^X$ and $g:2^Y \rightarrow \Re$ is strongly submodular then $f:2^X \rightarrow \Re$ given by $f(A) = g(\Gamma(A))$ for all $A \subseteq X$ is strongly submodular. %Lemma 9.1
\item If $w_x \geq 0$ for all $x \in X$ then the function $f(A) = \sum_{a \in A} w_a$ is strongly submodular. %pg 74
\item Suppose $f:2^X \rightarrow \Re$ is zero on the empty set. Then $f$ is strongly submodular if and only if 
\begin{equation}
\sum_{B:A \subseteq B} (-1)^{|B \setminus A|} f(B) \leq 0 \label{eq:strong_moebius}
\end{equation}
for all $A \subsetneq X$.  %Note: including A = emptyset. This is Theorem 9.4 (iv). 
\item If $f(\emptyset) = 0$ then $f$ is strongly submodular if and only if there are subsets $A_1,\ldots,A_m$ and non-negative weights $w_1,\ldots,w_m$ such that $f(A) = \sum_{i:A \cap A_i \neq \emptyset} w_i$. 
\item If $w_x \geq 0$ for all $x \in X$ then the function $f:2^X \rightarrow \Re$ with $f(\emptyset) = 0$ and $f(A) = \max\{w_a:a \in A\}$ for nonempty $A \subseteq X$ is strongly submodular. 
\item If $\omega_x \in \Re$ for all $x \in X$ then the function $g:2^X \rightarrow \Re$ with $g(A) = \max\{\omega_a:a \in A\}$ for nonempty $A \subseteq X$ and 
$g(\emptyset) \leq  \min\{\omega_a:a \in X\}$  is strongly submodular. 
\end{enumerate}
\end{proposition}

 \begin{theorem} \label{thm:submodular_match}
 Let $X$ be a finite set and $f:2^X \rightarrow \Re$ a non-decreasing and  submodular function with $f(\emptyset) = 0$. For all  $Y \subseteq X$ there is a strongly submodular function $g:2^X \rightarrow \Re$ such that $g(A) \leq f(A)$ for all $A \subseteq X$ and $g(A) = f(A)$ when $A =Y$ or $|A| \leq 1$. 
 \end{theorem}

 \begin{proof}
 We prove this by induction on $|Y|$.  When $|Y| = 0$, define $g$ by $g(\emptyset) = 0$ and $g(A)  = \max\{f(\{a\}): a \in A\}$ for nonempty $A \subseteq X$. Then $g$ is strongly submodular (Proposition~\ref{prop:basic_strong} (ix)),  $g(A) \leq f(A)$ for all $A \subseteq X$ and  $g(A) = f(A)$ when $A =Y$ or $|A| \leq 1$. \\
  
 Suppose that $k \geq 1$ and the induction hypothesis holds for when $|Y| = k-1$. Select $y \in Y$ and define $f_y:2^{X \setminus \{y\}} \rightarrow \Re$ by $f_y(A) = f(A \cup \{y\}) - f(\{y\})$ for all $A \subseteq X \setminus \{y\}$. By the induction hypothesis there is strongly submodular $g_y:2^{X \setminus \{y\}} \rightarrow \Re$ such that $g_y(A) \leq f_y(A)$ for all $A \subseteq X \setminus \{y\}$ and $g_y(A) = f_y(A)$ when $A =  Y \setminus \{y\}$ or when $|A| \leq 1$. 
 
Define $r: 2^X \rightarrow \Re$ by $r(\emptyset) = 0$ and 
\[r(A) = \max \{ f(\{a\}) + f(\{y\}) - f(\{a,y\}) :  a \in A\}\]
for nonempty $A \subseteq X$, noting that $r(A) \leq f(\{y\})$ for all $A$ and $r(A) = f(\{y\})$ when $y \in A$. Since $f$ is submodular and $f(\emptyset) = 0$ we have $f(\{a\}) + f(\{y\}) - f(\{a,y\}) \geq 0$ for all $a$.  We define $g:2^X \rightarrow \Re$ by $g(\emptyset) = 0$
 \[g(A) = r(A) + g_y(A \setminus \{y\}).\]
 By Proposition~\ref{prop:basic_strong} (x), (v) and (ii), $g$ is strongly submodular.
  
We show by cases that $g(A) \leq f(A)$ for all $A \subseteq X$ and $g(A) = f(A)$ when $A =Y$ or $|A| \leq 1$.  
For all nonempty $A \subseteq X$ such that $y \not \in A$  we have
 \begin{align*}
 g(A) & = g_y(A) + r(A) \\
 & \leq f_y(A) + \big(f(\{a\}) + f(\{y\}) - f(\{a,y\}) \big)& \mbox{ for some $a \in A$}\\
  & = f(A \cup \{y\}) + f(\{a\})  - f(\{a,y\}) \\
  & \leq f(A)
  \end{align*}
  by the submodularity inequality applied to $A$ and $\{a,y\}$. 
  
  If $y \in A$ then
  \begin{align*}
  g(A) & = g_y(A \setminus \{y\}) + r(A) \\
  & \leq f_y(A \setminus \{y\}) + f(\{y\}) \\
  & = f(A).
  \end{align*}
  If $A = Y$ then 
  \begin{align*}
  g(A) & = g_y(Y \setminus \{y\}) + r(Y) \\
  & = f_y(Y \setminus \{y\}) + f(\{y\})\\
  & = f(Y).
  \end{align*}
  For all $a \neq y$,
   \[g(\{a\}) = g_y(\{a\}) + \left( f(\{a\}) + f(\{y\}) - f(\{a,y\})  \right) = f(\{a\})\]
  while
   \[g(\{y\}) = g_y(\emptyset) + f(\{y\}) = f(\{y\})\]
and
\[g(\emptyset) =  g_y(\emptyset) + r(\emptyset) = 0.\]  
 
  By induction, the result holds for all $Y$. 
  \end{proof}

\end{document}
